\documentclass[10pt,a4paper]{amsart}
\usepackage[margin=19mm]{geometry}
\usepackage{amsmath,amssymb,mathtools}
\usepackage[scr=rsfs,cal=euler]{mathalfa}
\usepackage{xfrac}
\usepackage[unicode,hidelinks,bookmarks=false]{hyperref}
\newcommand{\ie}{i.e.\ }
\newcommand{\cf}{cf.\ }
\newcommand{\resp}{respectively\ }
\newcommand{\Subsection}{\subsection}
\providecommand{\nobreakdash}{\nobreak\hbox{-}}
\setlength{\emergencystretch}{2em}
\multlinegap0pt
\usepackage{xcolor}
\usepackage{amssymb}
\usepackage{enumerate}
\usepackage{tikz}
\usepackage{mathtools}
\usepackage{dsfont}
\newtheorem{lemma}{Lemma}
\DeclareMathOperator{\Proj}{Proj}
\newcommand{\R}{\mathbb{R}}
\newcommand{\Z}{\mathbb{Z}}
\title{Uniformly recurrent subalgebras in finite von Neumann algebras}
\author{Tattwamasi Amrutam, Pierre Fima,, Yongle Jiang}
\date{Source: arXiv:2606.09673v2 (CC BY 4.0)}
\begin{document}
\maketitle

\noindent\emph{Source and licence.} Uniformly recurrent subalgebras in finite von Neumann algebras. Version arXiv:2606.09673v2 (CC BY 4.0), distributed by arXiv under the Creative Commons Attribution 4.0 International licence, http://creativecommons.org/licenses/by/4.0/, as recorded in the arXiv licence metadata for this exact version and shown on the abstract page https://arxiv.org/abs/2606.09673v2. Full licence notice: \texttt{licenses/arxiv-2606.09673-CC-BY-4.0.txt}.\par\smallskip

\noindent\textit{Editorial scope.} Counted unit: the article's lemma LemClosedLineProjections (source0.tex lines 631--635). The article's complete original proof of it is delivered with it (source0.tex lines 636--725, one \texttt{proof} environment of 90 lines). No mathematics is written, repaired or re-derived: the statement and the proof are the article's own lines, in the article's own notation, and no retained line is substituted or re-flowed.  No further source-local context is needed: every cross-reference of every retained line resolves inside the two delivered spans.  Nothing was omitted from any retained span, so the package carries no omission record.  No retained line contains a citation, so the wrapper carries no bibliography and no bibliography entry.  No retained line contains a figure environment, an image-inclusion command or drawing code, so no image is delivered and the package declares no asset dependency.  Editorial formatting only, in addition: an AMS \texttt{amsart} preamble that restates the article's own declarations and macros used by the retained lines, namely the environments \texttt{lemma} on separate counters, together with the standard packages those lines need.  The retained environments print in this document as Lemma 1 in order of appearance, whereas the article prints the same items with the numbering of its own full document.  The complete native source of the article as distributed in the arXiv e-print for this exact version is bundled unchanged as \texttt{sources/202416/source0.tex}.
\begin{lemma}
\label{LemClosedLineProjections}
There exists a closed topological embedding $\mathbb R\hookrightarrow \Proj({\rm L}^\infty([0,1]))$, $t\mapsto q_t$
for the \({\rm L}^2\)-topology. Moreover, the embedding may be chosen so that $\tau(q_t)=\frac12$ for all $t\in\R$.
\end{lemma}

\begin{proof}
Consider the compact metric space $Y:=\{0,1\}^{\mathbb Z}\times[0,1]$ with Borel probability measure
$\nu
        :=
       \mu\otimes \lambda$ where $\mu:=\bigotimes_{n\in\mathbb Z}
        \left(\frac12\delta_0+\frac12\delta_1\right)$ and $\lambda$ is the normalized Lebesgue measure on $[0,1]$. Since \((Y,\nu)\) is a standard non-atomic probability space, there is a
trace-preserving isomorphism ${\rm L}^\infty(Y)\cong {\rm L}^\infty([0,1])$ and it suffices to construct a closed topological embedding $\mathbb R\hookrightarrow \Proj({\rm L}^\infty(Y))$.
\medskip
\noindent Write a point of \(Y\) as \((\omega,u)\), where $\omega=(\omega_n)_{n\in\mathbb Z}\in\{0,1\}^{\mathbb Z}$ and $u\in[0,1]$. For each \(n\in\mathbb Z\), put $R_n:=\{\omega:\omega_n=1\}$. Then $\mu(R_n)=\frac12$. For \(t\in\mathbb R\), write uniquely $t=n+s$, $n\in\mathbb Z$ and  $s\in[0,1)$. Define
$A_t:=R_n\times(s,1]\sqcup R_{n+1}\times[0,s]$ and  $q_t:=1_{A_t}\in {\rm Proj}({\rm L}^\infty(Y))$. Note that $\nu(A_t)=(1-s)\nu(R_n)+s\nu(R_{n+1})=\frac12$. Hence $\tau(q_t)=\frac12$ $\forall t\in\mathbb R$ and it suffices to show that    $t\mapsto q_t$ is a closed topological embedding for the \({\rm L}^2\)-topology.
\medskip
\noindent Let us show that $t\mapsto q_t$ is continuous. Note that, for all Borel sets $A,B\subseteq Y$,
$$\|1_A-1_B\|_2^2=\nu(A\triangle B).$$
Fix $n\in\Z$ and let us show continuity on $[n,n+1)$. If \(t=n+s\) and \(t'=n+s'\) lie in the same interval
\([n,n+1)\), then \(A_t\) and \(A_{t'}\) differ only on the strip $\min(s,s')<u\leq \max(s,s')$. Therefore
$\Vert q_t-q_{t'}\Vert_2^2
        =
        \nu(A_t\triangle A_{t'})
        \leq |s-s'|$. It remains to check continuity at the integers. Let \(m\in\mathbb Z\). By
definition, $A_m=R_m\times(0,1]\cup R_{m+1}\times\{0\}$ hence, in ${\rm L}^\infty(Y)$ one has $q_m=1_{{R_m}\times[0,1]}$. For $t=m+s$ with $s\in [0,1)$. Since
$$A_t\setminus\left(R_m\times[0,1]\right)\subset R_{m+1}\times[0,s]\text{ and }\left(R_m\times[0,1]\right)\setminus A_t\subset R_m\times[0,s]$$
we have $ \nu(A_t\triangle R_m)\leq s$. Hence, if \(t\to m\) from the right, say \(t=m+s\) with \(s\to0^+\), then $\Vert q_t-q_m\Vert_2^2=\nu(A_t\triangle R_m)\leq s\to0$. If \(t\to m\) from the left, say \(t=(m-1)+s\) with \(s\to1^-\), then
$$A_t\setminus(R_m\times[0,1])\subset R_{m-1}\times(s,1]\text{ and }(R_m\times[0,1])\setminus A_t\subset R_m\times(s,1].$$
Therefore $\nu(A_t\triangle(R_m\times[0,1]))\leq 1-s\to0$. This proves continuity at \(m\).
\medskip
\noindent We now show that the map is injective. Let $t,t'\in\R$ with $t\neq t'$ and write $t=n+s$ and $t'=m+r,$
with \(n,m\in\mathbb Z\) and \(s,r\in[0,1)\). Define $c_t,c_{t'}:[0,1]\to\Z$ by
\[
        c_t(u)=
        \begin{cases}
        n+1, & u\leq s,\\
        n, & u>s,
        \end{cases}
        \qquad
        c_{t'}(u)=
        \begin{cases}
        m+1, & u\leq r,\\
        m, & u>r.
        \end{cases}
\]
Then $(\omega,u)\in A_t\Leftrightarrow\omega_{c_t(u)}=1$. Since \(t\neq t'\), the step functions \(c_t\) and \(c_{t'}\) differ on the Borel set $U:=\{u\in[0,1]:c_t(u)\neq c_{t'}(u)\}$
of positive Lebesgue measure. More precisely, writing
$t=n+s$ and $t'=m+r$, with $n,m\in\mathbb Z$ and $s,r\in[0,1)$, one has
$$
\lambda(U)=
\begin{cases}
|s-r|, & n=m,\\
1-\max\{s-r,0\}, & m=n+1,\\
1-\max\{r-s,0\}, & n=m+1,\\
1, & |n-m|\geq 2.
\end{cases}
$$
In every case, $\lambda(U)>0$. Note that, for every \(u\in U\), $ \mu(\{\omega:\omega_{c_t(u)}\neq\omega_{c_{t'}(u)}\})
        =
        \frac12$. Therefore, by Fubini,
\[
        \nu(A_t\triangle A_{t'})
        =
        \int_0^1
        \mu(\{\omega:\omega_{c_t(u)}\neq\omega_{c_{t'}(u)}\})\,d\lambda(u)
       =
        \frac12\lambda(U)>0.
\]
Thus \(q_t\neq q_{t'}\).
\medskip
\noindent Finally, we show that the image is closed and the inverse is continuous on the image. We will use the following Claim.
\medskip
\noindent\textbf{Claim.} \textit{If \((t_m)_m\) be an unbounded sequence in \(\mathbb R\) then the sequence $(q_{t_m})_m$is not Cauchy in \({\rm L}^2(Y)\).}
\medskip
\noindent\textit{Proof of the Claim.} After passing to a subsequence, we may write $t_m=n_m+s_m$, $n_m\in\mathbb Z$, $s_m\in[0,1)$, with the pairs $\{n_m,n_m+1\}$ pairwise disjoint. Therefore, for every \(u\in[0,1]\) and $m\neq l\in\Z$, $\mu(\{\omega:1_{A_{t_m}}(\omega,u)\neq 1_{A_{t_\ell}}(\omega,u)\})=\frac12$. By Fubini,
\[
        \nu(A_{t_m}\triangle A_{t_\ell})
        =
        \int_0^1
        \mu(\{\omega:1_{A_{t_m}}(\omega,u)\neq 1_{A_{t_\ell}}(\omega,u)\})
        \,d\lambda(u)
        =
        \frac12.
\]
Hence
\[
        \|q_{t_m}-q_{t_\ell}\|_2^2
        =
        \nu(A_{t_m}\triangle A_{t_\ell})
        =
        \frac12.
\]
Thus the subsequence \((q_{t_m})_m\) is not Cauchy.$\hfill\qed$
\medskip
\noindent Let us show that the image is closed in ${\rm L}^2(Y)$. Let $(t_m)_m$ be a sequence in $\R$ such that $(q_{t_m})$ converges in ${\rm L}^2(Y)$. By the Claim, $(t_m)_m$ is bounded. Passing to a subsequence, we may assume $t_m\to t\in\mathbb R$. By continuity, $q_{t_m}\to q_t$. This shows that the image is closed. Finally, we show that the inverse is continuous on the image. Assume that $q_{t_m}\to q_t$ then, by the Claim, the sequence $(t_m)_m$ is bounded, so it suffices to show that $t$ is the only accumulation point of the sequence $(t_m)_m$. Consider a subsequence such that $t_{m_k}\to t'$. By continuity, $q_{t_{m_k}}\to q_{t'}$ and by uniqueness of limits $q_{t'}=q_t$. By injectivity, we have $t'=t$. It completes the proof.\end{proof}

\end{document}
